% !PRINTED 327
\weberpage{327}
Ist der Punkt $\xi_3$ ein Wendepunkt und $x_1 = 0$ die Wendetangente, so hat $f$ die Gestalt

\weq{4}{f = x_1 \varphi + x_2^3 \psi,}
worin $\varphi$ eine Form $(n-1)^{\text{ten}}$, $\psi$ eine Form $(n-3)^{\text{ten}}$ Grades ist.

Umgekehrt ist, wenn $f$ diese Form hat, der Punkt $\xi_3$ ein Wendepunkt, und $x_1 = 0$ die Wendetangente.

Eine gerade Linie, die die Curve $f$ in zwei getrennten Punkten berührt, heisst eine \gesperrt{Doppeltangente}. Nehmen wir eine solche Doppeltangente zur Linie $x_3 = 0$ und die Berührungspunkte als die auf $x_3 = 0$ gelegenen Ecken des Coordinatendreiecks, so muss, wenn $x_3 = 0$ in $f = 0$ eingesetzt wird, eine Gleichung entstehen, die sowohl $x_1$ als $x_2$ zum quadratischen Factor hat. Die Function $f$ muss also die Gestalt haben:

\weq{4}{f = x_3 \varphi + x_1^2 x_2^2 \psi,} % [sic] number (4) reprinted; page also numbers preceding display (4)
worin $\varphi$ eine Form $(n-1)^{\text{ten}}$ Grades, $\psi$ eine Form $(n-4)^{\text{ten}}$ Grades ist.

Etwas allgemeiner können wir auch sagen: Wenn $x_3 = 0$ eine Doppeltangente ist, so kann $f$ so dargestellt werden:

\weq{5}{f = x_3 \varphi + \chi^2 \psi,}
worin $\chi = 0$ die Gleichung eines Kegelschnittes ist. Die Berührungspunkte der Doppeltangente sind die Schnittpunkte von $x_3 = 0$ und $\chi = 0$.

\begin{center}
\S.\,87.

\gesperrt{Fundamentale Covarianten einer ternären Form.}
\end{center}

Im \S.\,60 des ersten Bandes haben wir den Begriff der Invarianten und Covarianten einer Form kennen gelernt. Auch in der Geometrie spielen diese Formen eine grosse Rolle.

Wenn durch die lineare Substitution \S.\,85, (1) die ternäre Form $f(x_1, x_2, x_3)$, die wir abgekürzt auch mit $f(x)$ bezeichnen, in $\varphi(y)$ übergeht, und wenn wir mit $a$ die Coëfficienten von $f$, mit $b$ die Coëfficienten von $\varphi$ und mit $r$ die Substitutionsdeterminante bezeichnen, so heisst eine Form $\Phi(x, a)$, die sowohl in Bezug auf die $x$ wie in Bezug auf die $a$ homogen ist, eine \gesperrt{Covariante} von $f(x)$, wenn sie der Bedingung

\weq{1}{\Phi(y, b) = r^\lambda \Phi(x, a)}
%CONT>>
